\documentclass[12pt,a4paper]{article}
\usepackage{amsmath,amssymb}
\usepackage[margin=2.54cm]{geometry}

\title{明确重新表述辅助变量}
\author{}
\date{}

\usepackage[UTF8]{ctex}
\begin{document}
\maketitle

\section*{符号与数据}
让$I,J,N\in\mathbb{Z}_{>0}$分别表示预约次数、患者类型以及观察到的场景数量。定义如下：
\[
[K]:=\{1,\ldots,K\},
\qquad
[a;b]:=\{a,a+1,\ldots,b\}.
\]
对于每种患者类型$j\in[J]$，让$I_j\in\mathbb{Z}_{>0}$能够满足……
\[
\sum_{j\in[J]} I_j=I.
\]
设$T\in\mathbb{Z}_{>0}$为规划范围，设$\beta\in(0,1]$为CVaR参数，设$\varphi\ge 0$为Wasserstein半径参数，设$p\ge 1$为其他参数。
\[
\boldsymbol{\tau}=(\tau_1,\ldots,\tau_J)^\top\in\mathbb{R}^J
\]
这些应该是特定类型延迟容忍度的向量。对于每一个$l\in[I]$和$n\in[N]$，让…
\[
\boldsymbol{z}_l^n=(z_{l1}^n,\ldots,z_{lJ}^n)^\top\in\mathbb{R}^J
\]
作为观测到的场景数据。对于每个$i\in[I]$，进行记录
\[
\boldsymbol{x}_i=(x_{i1},\ldots,x_{iJ})^\top.
\]

可排序的集合是
\[
\mathcal{X}:=\left\{
\boldsymbol{x}\in\{0,1\}^{I\times J}:\;
\sum_{j\in[J]}x_{ij}=1\ \forall i\in[I],\quad
\sum_{i\in[I]}x_{ij}=I_j\ \forall j\in[J]
\right\},
\]
而预约时间可行的集合就是
\[
\mathcal{Y}:=\left\{
\boldsymbol{y}\in\mathbb{Z}^I:
0\le y_1\le y_2\le\cdots\le y_I\le T
\right\}.
\]

\section*{决策变量}
决策变量为$\boldsymbol{x}$、$\boldsymbol{y}$，以及连续的辅助变量。
\[
\eta_i\in\mathbb{R}
\qquad \forall i\in\{0\}\cup[2;I],
\]
\[
\gamma_i^n\in\mathbb{R}
\qquad \forall n\in[N],\ i\in\{0\}\cup[2;I].
\]

\section*{辅助变量的重新表述}
对于每一个$i\in[2;I]$，决策依赖型经验模型中所述的经验延迟-CVaR表达式，可以通过辅助变量$\eta_i$和$\gamma_i^n$来表示。
\[
\left(1-\frac{1}{\beta}\right)\eta_i-y_i
+\frac{1}{N\beta}\sum_{n\in[N]}\gamma_i^n,
\]
受…的约束
\[
\gamma_i^n\ge y_t+
\sum_{l=t}^{i-1}\boldsymbol{x}_l^\top\boldsymbol{z}_l^n,
\qquad \forall n\in[N],\ i\in[2;I],\ t\in[i-1],
\]
\[
\gamma_i^n\ge y_i,
\qquad \forall n\in[N],\ i\in[2;I],
\]
以及
\[
\gamma_i^n\ge\eta_i,
\qquad \forall n\in[N],\ i\in[2;I].
\]
经验性 overtime-CVaR目标如下所示：
\[
\left(1-\frac{1}{\beta}\right)\eta_0-T
+\frac{1}{N\beta}\sum_{n\in[N]}\gamma_0^n,
\]
受…的约束
\[
\gamma_0^n\ge T,
\qquad \forall n\in[N],
\]
\[
\gamma_0^n\ge y_t+
\sum_{l=t}^{I}\boldsymbol{x}_l^\top\boldsymbol{z}_l^n,
\qquad \forall n\in[N],\ t\in[I],
\]
以及
\[
\gamma_0^n\ge\eta_0,
\qquad \forall n\in[N].
\]

\section*{优化模型}
确定性有限维重构方法
\begin{alignat}{2}
\min_{\boldsymbol{x},\boldsymbol{y},\boldsymbol{\eta},\boldsymbol{\gamma}}\quad
&\left(1-\frac{1}{\beta}\right)\eta_0-T
+\frac{1}{N\beta}\sum_{n\in[N]}\gamma_0^n
\label{eq:objective}\\[0.5em]
\text{s.t.}\quad
&\left(1-\frac{1}{\beta}\right)\eta_i-y_i
+\frac{1}{N\beta}\sum_{n\in[N]}\gamma_i^n
\le \boldsymbol{\tau}^{\top}\boldsymbol{x}_i
-\frac{\varphi}{\beta}(i-1)^{\frac{p-1}{p}},
&&\forall i\in[2;I],
\label{eq:delay-cvar}\\
&\gamma_i^n\ge y_t+
\sum_{l=t}^{i-1}\boldsymbol{x}_l^{\top}\boldsymbol{z}_l^n,
&&\forall n\in[N],\ i\in[2;I],\ t\in[i-1],
\label{eq:delay-scenario}\\
&\gamma_i^n\ge y_i,
&&\forall n\in[N],\ i\in[2;I],
\label{eq:appointment-scenario}\\
&\gamma_i^n\ge\eta_i,
&&\forall n\in[N],\ i\in[2;I]\cup\{0\},
\label{eq:threshold-scenario}\\
&\gamma_0^n\ge T,
&&\forall n\in[N],
\label{eq:overtime-horizon}\\
&\gamma_0^n\ge y_t+
\sum_{l=t}^{I}\boldsymbol{x}_l^{\top}\boldsymbol{z}_l^n,
&&\forall n\in[N],\ t\in[I],
\label{eq:overtime-scenario}\\
&\boldsymbol{x}\in\mathcal{X},\qquad \boldsymbol{y}\in\mathcal{Y}.&&
\label{eq:domains}
\end{alignat}

所有涉及$\boldsymbol{x}$以及场景向量$\boldsymbol{z}_l^n$和$\boldsymbol{\tau}^{\top}\boldsymbol{x}_i$的产品，在展开后都是仿射的，因为场景数据和容忍向量是固定的。辅助变量$\eta,\gamma\in\mathbb R$满足所显示的约束条件。

\end{document}
